\begin{ex}
Khối lập phương có bao nhiêu mặt phẳng?
\shortans{6}
\loigiai{
Khối lập phương có đúng 6 mặt phẳng bằng nhau.
\textbf{Đáp số:} $6$.
}
\end{ex}

\begin{ex}
Khối hộp chữ nhật có bao nhiêu mặt phẳng?
\shortans{6}
\loigiai{
Khối hộp chữ nhật có đúng 6 mặt phẳng.
\textbf{Đáp số:} $6$.
}
\end{ex}

\begin{ex}
Khối lập phương có bao nhiêu đỉnh?
\shortans{8}
\loigiai{
Khối lập phương có đúng 8 đỉnh.
\textbf{Đáp số:} $8$.
}
\end{ex}

\begin{ex}
Khối hộp chữ nhật có bao nhiêu cạnh?
\shortans{12}
\loigiai{
Khối hộp chữ nhật có đúng 12 cạnh.
\textbf{Đáp số:} $12$.
}
\end{ex}

\begin{ex}
Xếp 3 khối lập phương nhỏ thẳng hàng cạnh nhau. Khối mới tạo thành có dạng khối gì? (1: Khối lập phương; 2: Khối hộp chữ nhật).
\shortans{2}
\loigiai{
Khối mới có chiều dài gấp 3 lần chiều rộng nên là khối hộp chữ nhật (ghi số 2).
\textbf{Đáp số:} $2$.
}
\end{ex}

\begin{ex}
Bạn Hùng có 4 khối lập phương và 6 khối hộp chữ nhật. Hỏi bạn Hùng có tất cả bao nhiêu khối hình?
\shortans{10}
\loigiai{
Tổng số khối hình là: $4 + 6 = 10$ (khối hình).
\textbf{Đáp số:} $10$.
}
\end{ex}

\begin{ex}
Số khối hộp chữ nhật của Hùng nhiều hơn số khối lập phương bao nhiêu khối?
\shortans{2}
\loigiai{
Số khối hộp chữ nhật nhiều hơn số khối lập phương là: $6 - 4 = 2$ (khối).
\textbf{Đáp số:} $2$.
}
\end{ex}

\begin{ex}
Một cột cờ thẳng đứng có 3 lá cờ màu đỏ, xanh, vàng treo thẳng hàng từ trên xuống. Hỏi lá cờ ở giữa là lá cờ màu gì? (Đỏ: 1, Xanh: 2, Vàng: 3).
\shortans{2}
\loigiai{
Lá cờ ở giữa là lá cờ màu xanh (ghi số 2).
\textbf{Đáp số:} $2$.
}
\end{ex}

\begin{ex}
Trong một hàng dọc có 9 bạn học sinh. Bạn Nam đứng ở vị trí thứ 5 tính từ đầu hàng. Hỏi phía sau bạn Nam có bao nhiêu bạn học sinh?
\shortans{4}
\loigiai{
Số bạn học sinh phía sau Nam là: $9 - 5 = 4$ (bạn).
\textbf{Đáp số:} $4$.
}
\end{ex}

\begin{ex}
Trong một hàng dọc có 10 bạn học sinh. Bạn Lan đứng ở vị trí thứ 4 tính từ dưới lên. Hỏi phía trước bạn Lan có bao nhiêu bạn học sinh?
\shortans{6}
\loigiai{
Số bạn học sinh đứng trước Lan là: $10 - 4 = 6$ (bạn).
\textbf{Đáp số:} $6$.
}
\end{ex}

\begin{ex}
Một giá sách có 2 tầng: tầng trên có 5 khối đồ chơi hình lập phương, tầng dưới có 4 khối đồ chơi hình hộp chữ nhật. Em lấy bớt 3 khối ở tầng trên. Hỏi giá sách còn lại bao nhiêu khối đồ chơi cả hai tầng?
\shortans{6}
\loigiai{
Số khối đồ chơi còn lại ở tầng trên là: $5 - 3 = 2$ (khối).\nTổng số khối đồ chơi còn lại trên giá sách là: $2 + 4 = 6$ (khối đồ chơi).
\textbf{Đáp số:} $6$.
}
\end{ex}

\begin{ex}
Có 8 khối gỗ gồm khối lập phương và khối hộp chữ nhật. Biết số khối lập phương bằng số khối hộp chữ nhật. Hỏi có bao nhiêu khối lập phương?
\shortans{4}
\loigiai{
Vì hai loại có số lượng bằng nhau và tổng bằng 8, mà $4 + 4 = 8$ nên có $4$ khối lập phương.
\textbf{Đáp số:} $4$.
}
\end{ex}

\begin{ex}
Có 10 khối gỗ hình lập phương xếp thành một hàng dài. Bạn Mai lấy đi 2 khối ở đầu bên trái và 2 khối ở đầu bên phải. Hỏi hàng còn lại bao nhiêu khối gỗ?
\shortans{6}
\loigiai{
Tổng số khối gỗ Mai đã lấy đi là: $2 + 2 = 4$ (khối).\nSố khối gỗ còn lại trong hàng là: $10 - 4 = 6$ (khối gỗ).
\textbf{Đáp số:} $6$.
}
\end{ex}

\begin{ex}
Một tòa tháp gồm 3 tầng: tầng 1 (dưới cùng) có 4 khối gỗ, tầng 2 có 3 khối gỗ, tầng 3 (trên cùng) có 2 khối gỗ. Hỏi tòa tháp có tất cả bao nhiêu khối gỗ?
\shortans{9}
\loigiai{
Tổng số khối gỗ của tòa tháp là: $4 + 3 + 2 = 9$ (khối gỗ).
\textbf{Đáp số:} $9$.
}
\end{ex}

\begin{ex}
Tìm hiệu giữa số khối gỗ ở tầng 1 và tầng 3 của tòa tháp trên.
\shortans{2}
\loigiai{
Hiệu số khối gỗ giữa tầng 1 và tầng 3 là: $4 - 2 = 2$ (khối gỗ).
\textbf{Đáp số:} $2$.
}
\end{ex}

\begin{ex}
Một bể cá mini hình hộp chữ nhật có 4 mặt bên bằng kính và 1 mặt đáy bằng kính. Hỏi có tất cả bao nhiêu mặt kính bao quanh bể cá không có nắp?
\shortans{5}
\loigiai{
Bể cá không có nắp gồm 4 mặt bên và 1 mặt đáy: $4 + 1 = 5$ (mặt kính).
\textbf{Đáp số:} $5$.
}
\end{ex}

\begin{ex}
Xếp 4 khối lập phương nhỏ thành một hình vuông lớn nằm trên bàn. Hỏi có bao nhiêu khối lập phương nhỏ được sử dụng?
\shortans{4}
\loigiai{
Có đúng $4$ khối lập phương nhỏ được sử dụng.
\textbf{Đáp số:} $4$.
}
\end{ex}

\begin{ex}
Một chú kiến bò từ mép bàn bên trái sang mép bàn bên phải dài 10 gang tay. Chú kiến đã bò được 7 gang tay. Hỏi chú kiến còn phải bò bao nhiêu gang tay nữa để tới mép bàn bên phải?
\shortans{3}
\loigiai{
Quãng đường chú kiến còn phải bò là: $10 - 7 = 3$ (gang tay).
\textbf{Đáp số:} $3$.
}
\end{ex}

\begin{ex}
Trong một bức tranh: chú chim bay ở phía trên cây xanh, chú thỏ chạy ở phía dưới gốc cây. Hỏi có bao nhiêu con vật trong bức tranh?
\shortans{2}
\loigiai{
Có 2 con vật là chú chim và chú thỏ.
\textbf{Đáp số:} $2$.
}
\end{ex}

\begin{ex}
Một ngăn kéo bàn hình hộp chữ nhật chứa 7 quyển vở. Nam lấy ra 3 quyển vở mang đi học. Hỏi trong ngăn kéo còn lại bao nhiêu quyển vở?
\shortans{4}
\loigiai{
Số quyển vở còn lại trong ngăn kéo là: $7 - 3 = 4$ (quyển vở).
\textbf{Đáp số:} $4$.
}
\end{ex}

\begin{ex}
Mặt trên của một khối lập phương có dạng hình gì? (1: Hình tròn; 2: Hình tam giác; 3: Hình vuông; 4: Hình chữ nhật).
\shortans{3}
\loigiai{
Mặt trên của khối lập phương là hình vuông (ghi số 3).
\textbf{Đáp số:} $3$.
}
\end{ex}

\begin{ex}
Một chiếc hộp bánh hình lập phương có 6 mặt, người ta dán nhãn trang trí vào 4 mặt xung quanh. Hỏi còn bao nhiêu mặt chưa dán nhãn?
\shortans{2}
\loigiai{
Số mặt chưa dán nhãn là: $6 - 4 = 2$ (mặt).
\textbf{Đáp số:} $2$.
}
\end{ex}

\begin{ex}
Bên trái của bạn Bình có 3 bạn học sinh, bên phải của bạn Bình có 4 bạn học sinh. Hỏi hàng đó có tất cả bao nhiêu bạn học sinh (kể cả Bình)?
\shortans{8}
\loigiai{
Tổng số bạn trong hàng là: $3 + 1 + 4 = 8$ (bạn học sinh).
\textbf{Đáp số:} $8$.
}
\end{ex}

\begin{ex}
Trước mặt bạn Nam có 5 bạn học sinh, sau lưng bạn Nam có 4 bạn học sinh. Hỏi hàng dọc đó có tất cả bao nhiêu bạn học sinh (kể cả Nam)?
\shortans{10}
\loigiai{
Tổng số bạn trong hàng dọc là: $5 + 1 + 4 = 10$ (bạn học sinh).
\textbf{Đáp số:} $10$.
}
\end{ex}

\begin{ex}
Một người xếp 8 khối lập phương thành một khối lập phương lớn hơn. Sau đó nhấc bớt đi 1 khối ở góc trên. Hỏi còn lại bao nhiêu khối lập phương?
\shortans{7}
\loigiai{
Số khối lập phương còn lại là: $8 - 1 = 7$ (khối).
\textbf{Đáp số:} $7$.
}
\end{ex}
